\documentclass[11pt]{article}
\pagestyle{plain}

% \pgfplotsset{compat=1.18}
\usepackage[margin=1in]{geometry}
\usepackage{amsmath,amsfonts,amssymb,amsthm}
\usepackage[ruled,vlined,linesnumbered,commentsnumbered]{algorithm2e}
\SetKwInput{KwIn}{Input}
\SetKwInput{KwOut}{Output}
\SetKwInput{KwGlobal}{Global variable}
\SetKwInput{KwOracle}{Oracle}
\SetKwComment{Comment}{// }{}
\usepackage{graphicx}
\usepackage{enumerate}
\usepackage{bbm}
\usepackage{verbatim}
\usepackage{hyperref,color}
\usepackage[capitalize,nameinlink]{cleveref}
\usepackage[dvipsnames]{xcolor}
\hypersetup{
	colorlinks=true,
	pdfpagemode=UseNone,
	citecolor=OliveGreen,
	linkcolor=NavyBlue,
	urlcolor=Magenta,
	pdfstartview=FitW
}
\usepackage{appendix}
\usepackage{makecell}
\usepackage{tablefootnote}
\crefname{appsec}{Appendix}{Appendices}
\usepackage{tikz}
\usepackage{pgfplots}
\usepackage{xifthen}
\usepackage{subcaption}
\usepackage{hhline}
\usepackage{listings}

\theoremstyle{plain}
\newtheorem{theorem}{Theorem}[section]
\newtheorem{proposition}[theorem]{Proposition}
\newtheorem{lemma}[theorem]{Lemma}
\newtheorem{corollary}[theorem]{Corollary}
\newtheorem{conjecture}[theorem]{Conjecture}
\newtheorem{problem}[theorem]{Problem}
\newtheorem{claim}[theorem]{Claim}
\newtheorem{fact}[theorem]{Fact}

\theoremstyle{definition}
\newtheorem{definition}[theorem]{Definition}
\newtheorem{example}[theorem]{Example}
\newtheorem{observation}[theorem]{Observation}
\newtheorem{assumption}[theorem]{Assumption}
\newtheorem*{assumption*}{Assumption}
\newtheorem{condition}[theorem]{Condition}

\theoremstyle{remark}
\newtheorem{remark}[theorem]{Remark}

\crefname{lemma}{Lemma}{Lemmas}
\crefname{theorem}{Theorem}{Theorems}
\crefname{definition}{Definition}{Definitions}
\crefname{fact}{Fact}{Facts}
\crefname{claim}{Claim}{Claims}
\crefname{proposition}{Proposition}{Propositions}


\newcommand{\dif}{\,\mathrm{d}}
%\newcommand{\E}{\mathbb{E}}
%\newcommand{\Var}{\mathrm{Var}}
\newcommand{\Cov}{\mathrm{Cov}}
\newcommand{\MI}{\mathrm{I}}
%\newcommand{\Ent}{\mathrm{Ent}}
\newcommand{\ind}{\mathbbm{1}}
\newcommand{\allone}{\mathbf{1}}
\newcommand{\tr}{\mathrm{tr}}
\newcommand{\wt}{\mathrm{weight}}
\DeclareMathOperator*{\argmax}{arg\,max}
\newcommand{\norm}[1]{\left\lVert #1 \right\rVert}
\newcommand{\TV}[2]{\left\|#1 - #2\right\|_{\mathrm{TV}}}
\newcommand{\ceil}[1]{\left\lceil #1 \right\rceil}
\newcommand{\floor}[1]{\left\lfloor #1 \right\rfloor}
\newcommand{\ip}[2]{\left\langle #1 , #2 \right\rangle}
\newcommand{\grad}{\nabla}
\newcommand{\hessian}{\nabla^2}
\newcommand{\trans}{\intercal}
\newcommand{\diag}{\mathrm{diag}}
\newcommand{\poly}{\mathrm{poly}}
\newcommand{\dist}{\mathrm{dist}}
\newcommand{\sgn}{\mathrm{sgn}}
\newcommand{\fpras}{\mathsf{FPRAS}}
\newcommand{\fpaus}{\mathsf{FPAUS}}
\newcommand{\fptas}{\mathsf{FPTAS}}

\newcommand{\eps}{\varepsilon}

\newcommand{\N}{\mathbb{N}}
\newcommand{\Z}{\mathbb{Z}}
\newcommand{\Q}{\mathbb{Q}}
\newcommand{\R}{\mathbb{R}}
\newcommand{\F}{\mathbb{F}}
\newcommand{\C}{\mathbb{C}}

\renewcommand{\AA}{\mathcal{A}}
\newcommand{\BB}{\mathcal{B}}
\newcommand{\CC}{\mathcal{C}}
\newcommand{\DD}{\mathcal{D}}
\newcommand{\EE}{\mathcal{E}}
\newcommand{\FF}{\mathcal{F}}
\newcommand{\GG}{\mathcal{G}}
\newcommand{\HH}{\mathcal{H}}
\newcommand{\II}{\mathcal{I}}
\newcommand{\JJ}{\mathcal{J}}
\newcommand{\KK}{\mathcal{K}}
\newcommand{\LL}{\mathcal{L}}
\newcommand{\MM}{\mathcal{M}}
\newcommand{\NN}{\mathcal{N}}
\newcommand{\OO}{\mathcal{O}}
\newcommand{\PP}{\mathcal{P}}
\newcommand{\QQ}{\mathcal{Q}}
\newcommand{\RR}{\mathcal{R}}
\renewcommand{\SS}{\mathcal{S}}
\newcommand{\TT}{\mathcal{T}}
\newcommand{\UU}{\mathcal{U}}
\newcommand{\VV}{\mathcal{V}}
\newcommand{\WW}{\mathcal{W}}
\newcommand{\XX}{\mathcal{X}}
\newcommand{\YY}{\mathcal{Y}}
\newcommand{\ZZ}{\mathcal{Z}}

\newcommand{\KL}[2]{D_{\mathrm{KL}} \left( #1 \,\Vert\, #2 \right)}
\newcommand{\KLbig}[2]{D_{\mathrm{KL}} \left( #1 \,\big\Vert\, #2 \right)}
\newcommand{\KLBig}[2]{D_{\mathrm{KL}} \left( #1 \,\Big\Vert\, #2 \right)}

\newcommand{\SKL}[2]{D_{\mathrm{SKL}} \left( #1 \,\Vert\, #2 \right)}
\newcommand{\SKLbig}[2]{D_{\mathrm{SKL}} \left( #1 \,\big\Vert\, #2 \right)}
\newcommand{\SKLBig}[2]{D_{\mathrm{SKL}} \left( #1 \,\Big\Vert\, #2 \right)}

\newcommand{\qandq}{\quad\text{and}\quad}
\newcommand{\supp}{\mathrm{supp}}

\newcommand{\DKL}[2]{\-D_{\-{KL}}\left(#1 \parallel #2\right)}
\newcommand{\DTV}[2]{\-D_{\mathrm{TV}}\left({#1},{#2}\right)}
% \newcommand{\dist}{\mathrm{dist}}
\newcommand{\e}{\mathrm{e}}
\renewcommand{\epsilon}{\varepsilon}
\renewcommand{\emptyset}{\varnothing}
% \newcommand{\norm}[1]{\left\Vert#1\right\Vert}
\newcommand{\set}[1]{\left\{#1\right\}}
\newcommand{\tuple}[1]{\left(#1\right)} 
% \newcommand{\eps}{\varepsilon}
\newcommand{\inner}[2]{\left\langle #1,#2\right\rangle}
\newcommand{\tp}{\tuple}
\newcommand{\ol}{\overline}
\newcommand{\abs}[1]{\left\vert#1\right\vert}
\newcommand{\ctp}[1]{\left\lceil#1\right\rceil}
\newcommand{\ftp}[1]{\left\lfloor#1\right\rfloor}
\newcommand{\Ind}[1]{\-{Ind}_{#1}}
\newcommand{\todo}[1]{{\color{red} TODO: #1}}

\newcommand{\Cor}{\Psi^{\mathrm{sym}}}
\newcommand{\cor}{\Psi^{\-{cor}}}
\newcommand{\Inf}{\Psi}
% \newcommand{\diag}{\-{diag}}

\def\*#1{\boldsymbol{#1}} % Use \*A for \mathbf{A}
\def\+#1{\mathcal{#1}} % Use \+A for \mathcal{A}
\def\-#1{\mathrm{#1}} % Use \-A for \mathrm{A}
\def\=#1{\mathbb{#1}} % Use \^A for \mathbb{A}
\def\!#1{\mathfrak{#1}} % Use \!A for \mathfrak{A}

\newcommand*\diff{\mathop{}\!\mathrm{d}}

% probability
% \DeclareMathOperator*{\oPr}{\mathbf{Pr}}
\def\oPr{\mathop{\mathrm{Pr}}}
\renewcommand{\Pr}[2][]{ \ifthenelse{\isempty{#1}}
  {\oPr\left[#2\right]}
  {\oPr_{#1}\left[#2\right]} } % Use \Pr[a]{b} for \mathbf{Pr}_a[b], \Pr{b} for  \mathbf{Pr}[b]

% expectation: relies on the package "dsfont"
% \DeclareMathOperator*{\oE}{\mathbb{E}}
\def\oE{\mathop{\mathbb{E}}}
\newcommand{\E}[2][]{ \ifthenelse{\isempty{#1}}
  {\oE\left[#2\right]}
  {\oE_{#1}\left[#2\right]} }

% variance
%\DeclareMathOperator*{\oVar}{\mathbf{Var}}
\def\oVar{\mathrm{Var}}
\newcommand{\Var}[2][]{ \ifthenelse{\isempty{#1}}
  {\oVar\left[#2\right]}
  {\oVar_{#1}\left[#2\right]} }

% entropy
% \DeclareMathOperator*{\oEnt}{\mathbf{Ent}}
\def\oEnt{\mathrm{Ent}}
\newcommand{\Ent}[2][]{ \ifthenelse{\isempty{#1}}
  {\oEnt\left[#2\right]}
  {\oEnt_{#1}\left[#2\right]} }

\newcommand{\PhiEnt}[2][]{ \ifthenelse{\isempty{#1}}
  {\oEnt^\phi\left[#2\right]}
  {\oEnt^\phi_{#1}\left[#2\right]} }


\newenvironment{Exampleparagraph}[1][]
{
    \par 
    \noindent 
    \textbf{\color{blue} #1}
    \par 
    \vspace{0.5em} 
    \itshape 
}
{
    \par 
    \vspace{1em} 
}

%\input{local-macro}

\newcommand{\RED}[1]{\textcolor{red}{#1}}

\newcommand{\bern}[1]{\mathrm{Bern}\left(#1\right)}
\newcommand{\pll}{\parallel}



\title{MATH 6112. FALL 2026. ASSIGNMENT 2.}
\date{Due on October 16, 2026}
\author{{\textbf{Zejia Chen}} \\ \textbf{GTID: 904178159}}

\pgfplotsset{compat=1.18} 
\begin{document}
\maketitle
Let $D_p = \begin{pmatrix} I_p & 0\\0 & 0\end{pmatrix}$ and we use $e_i$ to denote the $i$-th standard basis vector of $\C^n$, with $1$ in the $i$-th coordinate and $0$ elsewhere.
\section*{Problem 1}
\textbf{(a)}
$\boldsymbol{\implies}$: for any eigenvalue $\lambda$ of $B$ with eigenvector $v$, we have that $B^2v = \lambda^2 v = Bv = \lambda v$.
Since $v\neq 0$, $\lambda^2 - \lambda = 0 \implies \lambda = 0,1$.
Therefore, diagonalization of $B$ implies it is similar to $D_p$ for some $p\in\{0\}\cup[n]$.
It is sufficient to show that $B$ is diagonalizable, that is $\dim \ker B + \dim \ker (I-B) = n$.
Suppose $\-{rank}(B) = n-k$, then $\dim \ker B = k$.
Since $\dim \ker (I-B)+\dim \-{Im} (I-B) = n$, we need to show $\dim \-{Im} (I-B) = k$.
Since for any $x\in \C^n$, $x = Bx + (I-B)x$ and hence $\-{span}(\-{Im} B , \-{Im}(I-B)) = \C^n$.
For $y\in \-{Im} B \cap \-{Im}(I-B)$, let $y = Bt_1 = (I-B)t_2$, we have that
$$
    0 = (B-B^2)t_2 = By = B^2 t_1 = Bt_1 = y .
$$
Therefore, $\-{Im} B \cap \-{Im}(I-B) = \{0\}$ and since $\-{span}(\-{Im} B, \-{Im}(I-B)) = \C^n$, $\dim \-{Im} B + \dim \-{Im}(I-B) = n$.
We know $\dim \-{Im} B = \-{rank}(B) = n-k$, thus, $\dim \-{Im}(I-B) = k$.
In conclusion, $B$ is diagonalizable.
\\$\boldsymbol{\impliedby}$: similarity implies that there exists $V\in \C^{n\times n}$ such that $B = V^{-1} D_p V$.
Since $D_p^2=D_p$, $B^2 = V^{-1} D_p V V^{-1} D_p V = V^{-1} D_p V = B$.
\\\textbf{(b)} Suppose $V = [v_1,\cdots,v_n]$.
Since $V^{-1}BV = D_p$, we have that $Bv_i = v_i$ for $i\in[p]$ and $Bv_i = 0$ for $i > p$.
For any vector $v \in \C^n$, suppose $v = \sum_{i=1}^n \lambda_i v_i$ since $\-{span}(V) = \C^n$ and we have that
$$
    Bv = \sum_{i=1}^n \lambda_i Bv_i = \sum_{i=1}^p \lambda_i Bv_i \in \-{span}\{v_1,\ldots,v_p\}.
$$ 
Together with $Bv_i = v_i$ for $i\in[p]$, $B$ projects onto the subspace $\+V = \-{span}\{v_1,\ldots,v_p\}$.
\\\textbf{(c)}
Note that the plane $2x_1 - 3x_2 + x_3 = 0$ is 
$$
    \-{span}\left\{\begin{pmatrix}
    0\\1\\3
    \end{pmatrix},\begin{pmatrix}
    -1\\0\\2
    \end{pmatrix}\right\}. 
$$
Define 
$$
    V = \begin{pmatrix}
           0 & -1 & 0
        \\ 1 & 0 & 1
        \\ 3 & 2 & 0
    \end{pmatrix}\quad \text{and}\quad B = V D_2 V^{-1} = \frac 13\begin{pmatrix}
           3 & 0 & 0
        \\ 2 & 0 & 1
        \\ 0 & 0 & 3
    \end{pmatrix} .
$$
By Part \textbf{(b)}, $B$ projects onto the subspace $2x_1-3x_2+x_3=0$ and $B^* \neq B$.
Take Gram-Schmidt Orthogonalization of $V$, we have that 
$$
    U = \frac 1{\sqrt{70}}\begin{pmatrix}
           0 & -5\sqrt{2} & -2\sqrt{5}
        \\ \sqrt 7 & -3\sqrt 2 & 3\sqrt 5
        \\ 3\sqrt 7 & \sqrt 2 & -\sqrt 5
    \end{pmatrix} \quad \text{and} \quad \-{span}(U_1,U_2) = \{x\in \R^{3}\mid 2x_1-3x_2+x_3=0\},
$$
where we use $U_i$ to denote the $i$-th column of $U$.
Take $B' = UD_2U^* = \frac 1{14}\begin{pmatrix}10 & 6 & -2\\6 & 5 & 3\\-2 & 3 & 13\end{pmatrix}$.
By Part \textbf{(b)}, $B'$ projects onto the subspace $2x_1-3x_2+x_3=0$ and $(B')^* = B'$.
\\\textbf{(d)} 
Take $V = \begin{pmatrix}1 & t\\0 & 1\end{pmatrix}$ and $B := VD_1V^{-1} =  \begin{pmatrix}1 & -t\\0 & 0\end{pmatrix}$.
By Part \textbf{(b)}, $B$ projects onto $\-{span}\{e_1\}$ and $B^\top \neq B$.
We have that $\norm{B}_2^2 \geq \norm{B(e_1-e_2)}_2^2/2 = (t+1)^2/2$ and take $t\to \infty$, we get arbitrarily large $2$-norm of $B$.
\section*{Problem 2}
\textbf{(a)} Since $A$ is normal, there exists a unitary transformation $U$ such that $A = U \Lambda U^*$ where $\Lambda = \diag\{\lambda_i\}_{i\in[n]}$ is the diagonal matrix of eigenvalues.
Let $U = [u_1,\cdots,u_n]$.
For $x\in \C^n$, suppose $x = \sum_{i=1}^n \mu_iu_i$ with $\sum_{i=1}^n\abs{\mu_i}^2 = 1$ and we have that 
$$
    x^*Ax = \sum_{i,j=1}^n \overline{\mu_i}\mu_j (u_i^*Au_j) = \sum_{i=1}^n \abs{\mu_i}^2 \lambda_i \in \-{Co}(\sigma(A)).
$$
Hence, $\+F(A) \subseteq \-{Co}(\sigma(A))$.
Fix $x =  \sum_{i=1}^n \gamma_i\lambda_i \in \-{Co(\sigma(A))}$ with $\sum_{i=1}^n \gamma_i=1$ and $\gamma_i \geq 0$ for $i\in[n]$, take $\mu_i = \sqrt{\gamma_i}$ and $y = \sum_{i=1}^n \mu_i u_i$.
We have that $\norm{y}_2=1$ and $y^*Ay = x$.
Therefore, $\-{Co}(\sigma(A)) \subseteq \+F(A)$.
In conclusion, $\-{Co}(\sigma(A)) = \+F(A)$.
\\\textbf{(b)} $\boldsymbol{\implies}$: $A$ is Hermitian implies all the eigenvalues of $A$ are real.
By Part \textbf{(a)}, $\+F(A)$ is a segment of the real line.
\\$\boldsymbol{\impliedby}$: for any $k\in [n]$, $e_k^\top Ae_k = A_{kk}\in \+F(A)\subseteq \R$.
Fix $k,j\in[n]$ with $k\neq j$.
We have that
$$
    \frac 12(e_k+e_j)^\top A(e_k+e_j) \in \+F(A) \subseteq \R \implies A_{kj} + A_{jk}\in \R,
$$
since $A_{kk},A_{jj}\in \R$.
Suppose $A_{kj} = \alpha_1 + \beta i$ and $A_{jk} = \alpha_2 - \beta i$ where $\alpha_1,\alpha_2,\beta\in \R$.
We have that 
$$
    \frac 12(e_k-ie_j)^\top A(e_k+ie_j) \in \+F(A) \subseteq \R \implies i(A_{kj} - A_{jk})\in \R \implies (\alpha_1-\alpha_2)i\in\R \implies \alpha_1=\alpha_2,
$$
which implies $A_{kj} = \overline{A_{jk}}$.
Therefore, $A$ is Hermitian.
\\The endpoints of $\+F(A)$ are $\lambda_{\min}(A)$ and $\lambda_{\max}(A)$. 
\\\textbf{(c)}
Define $P = \begin{pmatrix}
    0 & 1 \\ 0 & 0
\end{pmatrix}$.
For any $x\in \C^2$ with $\norm{x}_2=1$, we can write $x = (\cos\alpha e^{i\theta_1},\sin\alpha e^{i\theta_2})$, where $\alpha \in [0,\frac \pi2]$ and $\theta_1,\theta_2\in [0,2\pi)$, by polar transformation.
We have that $x^*Px = \overline{x_1}x_2 = \frac{1}{2}\sin(2\alpha)e^{i(\theta_2-\theta_1)}$.
Taking over $\alpha \in [0,\frac \pi2]$ and $\theta_1,\theta_2\in [0,2\pi)$, we conclude that $\+F(P)$ is a disk with radius $\frac 12$ and center $0$.

It is clear that $\+F(A) = \+F(P+I) = 1+\+F(P)$ and $r(A) \leq 1 + \frac 12$ and the equality can be achieved since $\+F(P)$ is a disk.
In the same way,  $\+F(A^2) = \+F(2P+I) = 1+2\+F(P)$ and $r(A^2) = 1+2\times \frac 12$.
Therefore, $r(A) = 3/2$ and $r(A^2) = 2$.
\\\textbf{(d)} Nonnegativity and positive homogeneity of numerical radius follow from the nonnegativity and positive homogeneity of the norm of complex number.
If $r(A) = 0$, then $\+F(A) = \{0\}$ and $A = 0$.
The triangle inequality follows from the following inequality:
\begin{align*}
    r(A+B) &= \max_{\norm{x}_2=1}\abs{x^*(A+B)x}
    \leq \max_{\norm{x}_2=1}\abs{x^*Ax}+\abs{x^*Bx}
    \\&\leq \max_{\norm{x}_2=1}\abs{x^*Ax}+\max_{\norm{y}_2=1}\abs{y^*Ay}
    = r(A)+r(B).
\end{align*}
Hence, $r(\cdot)$ is a norm.
However, consider $A = \begin{pmatrix} 0 & 1 \\ 0 & 0 \end{pmatrix}$ and $B = \begin{pmatrix} 0 & 0 \\ 1 & 0 \end{pmatrix}$.
We have that
$$
    r(AB) \geq |e_1^\top AB e_1| = 1.
$$
By AM-GM inequality,
$$
    \forall x\in \C^2 \text{ with }\norm{x}_2 = 1, \abs{\overline{x_1}x_2} = \abs{x_1}\abs{x_2} \leq \frac 12\tp{\abs{x_1}^2+\abs{x_2}^2} = \frac 12 \implies r(A) \leq \frac 12.
$$
In the same way, $r(B)\leq \frac 12$.
Therefore, $r(AB) = 1 > \frac 14 \geq r(A)r(B)$ and $r$ is not a consistent matrix norm.
\section*{Problem 3}
\begin{align*}
    \begin{pmatrix}
        A & 0
        \\I & B
    \end{pmatrix} &\xrightarrow{\text{swap columns}} 
    \begin{pmatrix}
        A & 0
        \\I & B
    \end{pmatrix}
    \begin{pmatrix}
        0 & I
        \\I & 0
    \end{pmatrix}
    =
    \begin{pmatrix}
        0 & A
        \\B & I
    \end{pmatrix}
    \\&\xrightarrow{\text{Add ($-A\times$row $2$) to row $1$}} 
    \begin{pmatrix}
        I & -A
        \\0 & I
    \end{pmatrix}
    \begin{pmatrix}
        0 & A
        \\B & I
    \end{pmatrix}
    =
    \begin{pmatrix}
        -AB & 0
        \\B & I
    \end{pmatrix} 
    \\&\xrightarrow{\text{Multiply row $1$ by $-1$}} 
    \begin{pmatrix}
        -I & 0
        \\0 & I
    \end{pmatrix} 
    \begin{pmatrix}
        -AB & 0
        \\B & I
    \end{pmatrix} 
    =
    \begin{pmatrix}
        AB & 0
        \\B & I
    \end{pmatrix}.
\end{align*}
Therefore,
$$
    M = \begin{pmatrix}
        -I & 0
        \\0 & I
    \end{pmatrix} \begin{pmatrix}
        I & -A
        \\0 & I
    \end{pmatrix}
    =\begin{pmatrix}
        -I & A
        \\0 & I
    \end{pmatrix}
    \quad \text{ and } \quad
    W = \begin{pmatrix}
        0 & I
        \\I & 0
    \end{pmatrix} .
$$
By properties in the problem statement, we claim that 
\begin{align*}
    \det(A)\det(B)&=\det \begin{pmatrix}
        A & 0
        \\I & B
    \end{pmatrix} =
    (-1)^n\det \begin{pmatrix}
        0 & A
        \\B & I
    \end{pmatrix}
    = 
    (-1)^n \det \begin{pmatrix}
        -AB & 0
        \\B & I
    \end{pmatrix}
    \\&= 
    (-1)^{2n}\det \begin{pmatrix}
        AB & 0
        \\B & I
    \end{pmatrix}=\det(AB) .
\end{align*}
\section*{Problem 4}
We first show that for any matrix $A\in \F^{n\times n}$, $\norm{A^*}_2 = \norm{A}_2$:
\begin{align*}
    \norm{A}_2 = \sup_{\norm{x}_2 = 1} \norm{Ax}_2
    = \sup_{\norm{x}_2 = \norm{y}_2 = 1} \abs{\inner{Ax}{y}}
    = \sup_{\norm{x}_2 = \norm{y}_2 = 1} \abs{\inner{x}{A^*y}}
    = \sup_{\norm{y}_2 = 1} \norm{A^*y}_2 = \norm{A^*}_2 .
\end{align*}
Therefore, we have that $\norm{(A^*)^{-1}}_2\norm{A-B}_2 < 1$.
Suppose $B$ is not invertible, that is, there is a vector $v\in \F^n$ such that $\norm{v}_2 = 1$ and $Bv = 0$.
Hence, we have that 
\begin{align*}
\norm{(A^*)^{-1}}_2\norm{A-B}_2 &\geq \norm{(A^*)^{-1}v}_2*\norm{(A-B)v}_2 = \norm{(A^*)^{-1}v}_2*\norm{Av}_2
\\& \geq \abs{\inner{(A^*)^{-1}v}{Av}} = \abs{v^*(A^{-1}A)v} = 1,
\end{align*}
which contradicts $\norm{(A^*)^{-1}}_2\norm{A-B}_2 < 1$.
\section*{Problem 5}
Simple observation shows that
$$
    \begin{pmatrix}
        A & B
        \\ C & D
    \end{pmatrix} = 
    \begin{pmatrix}
        I & 0
        \\ CA^{-1} & D-CA^{-1}B
    \end{pmatrix}
    \begin{pmatrix}
        I & -B
        \\ 0 & I
    \end{pmatrix}^{-1}
    \begin{pmatrix}
        A^{-1} & 0
        \\ 0 & I
    \end{pmatrix}^{-1} .
$$
Therefore, by Problem \textbf{(3)}, we have that
\begin{align*}
    \det(M) &= \det\begin{pmatrix}
        I & 0
        \\ CA^{-1} & D-CA^{-1}B
    \end{pmatrix}
    \det\begin{pmatrix}
        I & -B
        \\ 0 & I
    \end{pmatrix}^{-1}
    \det \begin{pmatrix}
        A^{-1} & 0
        \\ 0 & I
    \end{pmatrix}^{-1}
    \\&= \det(D-CA^{-1}B)\det(A^{-1})^{-1}
    \\&= \det(A)\det(D-CA^{-1}B)
    =\det(AD-ACA^{-1}B) &&(\textbf{Problem }(3))
    \\&=\det(AD-CAA^{-1}B) = \det(AD-CB), &&(AC=CA)
\end{align*}
where the second equality follows from $\det(A^{-1}) = \det(A)^{-1}$ and the determinant of a block matrix. 

For the counterexample, let us consider $n = 2$ and
$$
    M = \begin{pmatrix}
          1& 0& 1& 0
        \\1& 2& 0& 1
        \\1& 1& 1& 0
        \\1& 1& 0& 2
    \end{pmatrix}.
$$
We have that $\det(M) = 1$ while $\det(AD-CB) = 0$.

Suppose $A$ is not invertible.  
The eigenvalues of $A$ are $\lambda_1,\cdots,\lambda_n$ and let
$$
    m = \begin{cases}1 &,\lambda_1 = \cdots=\lambda_n=0 \\\min_{i:\lambda_i\neq 0}|\lambda_i| &,\text{otherwise} \end{cases} .
$$
Take $t\in (0,m)$, we have that $A + t I$ is invertible: 
$$
    \exists V\in C^{n\times n}: A=V^{-1}\tp{R+\diag\{\lambda_i\}_{i\in[n]}}V \implies A+tI = V^{-1}\tp{R+\diag\{\lambda_i+t\}_{i\in[n]}}V,
$$
where $R$ is some strict upper triangular matrix and $\lambda_i + t \neq 0$ for $i\in[n]$.
And also, $(A+tI)C = AC+tC=CA+tC=C(A+tI)$.
Then by the previous part, we have that 
$$
    \det\begin{pmatrix}
        A + tI & B
        \\ C & D
    \end{pmatrix} = \det(AD+tD-CB).
$$
Note that $\det\begin{pmatrix}
        A + tI & B
        \\ C & D
    \end{pmatrix}$ is a polynomial function with respect to $t$, and thereby continuous.
Taking $t$ tends to $0$, we conclude that
$$
    \det\begin{pmatrix}
        A & B
        \\ C & D
    \end{pmatrix} = \lim_{t\to 0^+}\det\begin{pmatrix}
        A+tI & B
        \\ C & D
    \end{pmatrix} = \lim_{t\to 0^+}\det(AD+tD-CB) = \det(AD-CB).
$$
\section*{Problem 6}
\textbf{(a)} Diagonalization of $A$ holds since $A$ is Hermitian.
Let $U,V$ be unitary matrices satisfying that
$$
    U^*AU=V^*AV = \begin{pmatrix}
        \mu_1 I_{n_1} & & 
        \\& \ddots &
        \\& & \mu_pI_{n_p}
    \end{pmatrix}\quad\text{ and }\quad
    W = U^*V = \begin{pmatrix}
        W_{11} &\cdots & W_{1p}
        \\\vdots& \ddots &\vdots
        \\W_{p1}&\cdots & W_{pp}
    \end{pmatrix}.
$$
Let $D = \diag\{\mu_i I_{n_i}\}_{i\in[p]}$.
Following from the above equality, we have that $WD=DW$, that is,
$$
    \begin{pmatrix}
        \mu_1W_{11} &\cdots & \mu_1W_{1p}
        \\\vdots& \ddots &\vdots
        \\\mu_pW_{p1}&\cdots & \mu_pW_{pp}
    \end{pmatrix} = \begin{pmatrix}
        \mu_1W_{11} &\cdots & \mu_pW_{1p}
        \\\vdots& \ddots &\vdots
        \\\mu_1W_{p1}&\cdots & \mu_pW_{pp}
    \end{pmatrix} .
$$
For $i\neq j$, $\mu_i\neq \mu_j$ but $\mu_iW_{ij}=\mu_jW_{ij}$, implying $W_{ij}=0$.
Therefore $W = \diag\{W_{ii}\}_{i\in[p]}$.
Since $\diag\{W_{ii}^*W_{ii}\}_{i\in[p]} = W^*W = V^*UU^*V = I$, we have that $W_{ii}^*W_{ii}=I_{n_i}$ for any $i\in[p]$.
This completes the proof.
\\\textbf{(b)} The proof is exactly the same. 
Diagonalization of $A$ holds since $A$ is symmetric.
Let $U,V$ be orthogonal transformations and $U^\top A U = V^\top A V = \diag\{\mu_iI_{n_i}\}_{i\in[p]}$.
Let $D = \diag\{\mu_i I_{n_i}\}_{i\in[p]}$ and $W=U^\top V$.
We have that $DW=WD$, implying $W = \diag\{W_{ii}\}_{i\in[p]}$ is diagonal block matrix.
Since $W = U^\top V$, $W$ is orthogonal and thus, $W_{ii}$ is orthogonal for $i\in [p]$.
This completes the proof.
% \section*{Acknowledgement}
% External tools, such as AI and WolframAlpha, were used to calculate determinants and matrix multiplication, and provide some hints.
\end{document}